El moviment d'una anguila es pot aproximar al d'una ona harmònica transversal que es propaga des del cap fins a la cua. Per a estudiar-ne el moviment hem de simplificar: no considerem l'aportació al moviment de la resta de músculs del cos, i suposem, simplement, que l'ona es genera al cap de l'anguila, que vibra amb una freqüència de 0,50 Hz i amb una amplitud de 5,00 cm. La distància entre dos punts consecutius del cos de l'anguila que estan en el mateix estat de vibració és de 20,0 cm.

\figura[0.6]{fig1.png}

\begin{apartat}{1,25}
Calculeu la velocitat a la qual es propaga l'ona pel cos de l'anguila, la freqüència angular i el nombre d'ona. Si a l'instant inicial el cap té una elongació zero i la velocitat d'oscil·lació transversal és positiva, determineu l'expressió de l'equació d'ona. Per a l'equació d'ona utilitzeu el sistema de coordenades de la figura superior, on el cap de l'anguila es troba sempre a l'origen de les abscisses.
\end{apartat}

\begin{apartat}{1,25}
A partir de l'equació d'ona, deduïu i calculeu els mòduls de la velocitat i de l'acceleració màximes de l'oscil·lació transversal. Si la longitud total de l'anguila és de 58,0 cm, calculeu també la velocitat i l'acceleració transversals a la cua 10 s més tard d'haver iniciat el moviment.
\end{apartat}
